\documentclass[12pt]{article}
\usepackage[T2A]{fontenc}
\usepackage[utf8]{inputenc}
\usepackage[russian]{babel}
\usepackage{latexsym,amssymb,amsthm}
\usepackage{amsmath}
\usepackage{wrapfig}
\RequirePackage[pdftex]{graphicx}
%%\usepackage{array}

\usepackage{epsf,epic,eepic}
\input ris.tex
% \input cells.sty
% \addrisoption{\risleftfalse}

\DeclareRobustCommand{\No}{\ifmmode{\nfss@text{\textnumero}}\else\textnumero\fi}

\oddsidemargin=0pt
\textwidth=159mm
\headheight=0pt
\headsep=0pt
\topmargin=0pt
\textheight=240mm

%%
%% Три точки для обозначения делимости
%%
\catcode`\@=11
\def\kratno{\mathrel{\smash{\lower.5ex\hbox{$\,\vdots\,$}}}}
% Это прямые три точки для обозначения делимости
\def\nekratno{\mathrel{\mathpalette\c@ncel\kratno}}
\def\c@ncel#1#2{\m@th\ooalign{$\hfil#1\mkern1mu/\hfil$\crcr$#1#2$}}
% Это перечеркнутые три точки для обозначения неделимости
\def\divis{\smash{\lower.1ex\hbox{\,\hbox{.}\kern-.22em\lower-.6ex\hbox{.}
\kern-.52em\lower-1.2ex\hbox{.}\,}} }
% А это наклонные...

\hyphenation{чис-ло чис-ла чис-лу чис-лом чис-ле чис-лам чис-ла-ми чис-лах
че-ты-рех-уголь-ник че-ты-рех-уголь-ни-ка пря-мо-уголь-ник тре-уголь-ни-ка
тре-уголь-ни-ков вне-впи-сан-ной вне-впи-сан-ных}

\sloppy
\pagestyle{empty}

\def\q#1.{{\bf #1. }}
\def\dotline{\smallskip\hbox to \hsize{\dotfill}\medskip}

\newbox\kk

\begin{document}

\centerline{\sc САНКТ-ПЕТЕРБУРГСКАЯ}
\centerline{\sc ОЛИМПИАДА ШКОЛЬНИКОВ 2025 года ПО МАТЕМАТИКЕ}
\centerline{\sc II тур. \ 9 класс.}
\medskip
\hrule
\bigskip

\q1. \emph{Полупростыми} будем называть натуральные числа, равные произведению двух различных простых чисел. Пусть $a<b<c<d$ --- четыре последовательных полупростых числа. Оказалось, что между числами $a$ и $b$ есть квадрат простого числа, а также между числами $c$ и~$d$ есть квадрат простого числа. Докажите, что между числами $b$ и $c$ нет квадратов простых чисел.

\q2. Дан выпуклый четырехугольник $ABCD$, в котором $AB = BC = CD \neq AD$. Биссектрисы углов $A$ и $C$ пересекаются внутри четырехугольника. Докажите, что биссектрисы углов $B$ и $D$ не могут пересечься внутри четырехугольника.

\q3. При каком максимальном $n$ можно 
отметить на шахматной доске $100\times 100$ некоторые поля так, чтобы нашлись два таких отмеченных поля $A$ и $B$, что шахматный король, двигаясь по отмеченным полям, мог бы пройти от $A$ до $B$ за $n$ ходов, но не мог бы пройти менее чем за $n$ ходов?

\setlength{\multlinegap}{0pt}

\q4. Положительные числа $a_1, a_2, \ldots, a_n$ таковы, что $a_1^2+\ldots+a_n^2 =1$. Докажите неравенство
\begin{multline*}
\frac{1}{a_1-a_1^3} + \frac{1}{a_2-a_2^3} + \ldots +\frac{1}{a_n-a_n^3} 
\geq\\ \geq 
\left(a_1+a_2+\ldots+a_n \right) \left( \frac{1}{1-a_1^2} + \frac{1}{1-a_2^2} + \ldots + \frac{1}{1-a_n^2} \right).
\end{multline*}

\dotline

\centerline{Олимпиада 2025 года. II тур. 9 класс. Выводная аудитория.}
\smallskip

\q5. На нескольких кустах растут ягоды. Каждый день с 1-го по 30-е июля ягоды вырастают по следующему правилу. Если $k$-го июля стоит хорошая погода, то ровно в полдень на всех кустах, где меньше $k$ ягод, вырастет по одной новой ягоде (на остальных кустах число ягод не меняется). Если же погода плохая, то в полдень по одной ягоде вырастет на всех кустах, где не меньше чем $k$ ягод. Вечером 30-го июля оказалось, что количества ягод на кустах различны и не превосходят 30. Найдите наибольшее возможное число кустов.

\q6. Натуральные числа $a$, $b$, $c$ и $n$ таковы, что $a^2-nab+b^2=c^2$. Докажите, что число $a+b+(n+1)c$ --- составное.

\q7. В треугольнике $ABC$ проведены медиана $AM$ и биссектриса $AL$. Высота, проведенная к стороне $BC$, пересекает вписанную окружность в точках $X$ и $Y$, а продолжение этой высоты пересекает одну из  вневписанных окружностей в точках $Z$ и $T$. Докажите, что $\angle XLY  = \angle ZMT$.


\clearpage

\centerline{\sc САНКТ-ПЕТЕРБУРГСКАЯ}
\centerline{\sc ОЛИМПИАДА ШКОЛЬНИКОВ 2025 года ПО МАТЕМАТИКЕ}
\centerline{\sc II тур. \ 10 класс.}
\medskip
\hrule
\medskip

\q1. В ряд лежат 39 одинаковых с виду монет. Какие-то 10 подряд лежащих
монет~--- фальшивые. Все настоящие монеты весят одинаково, все фальшивые
тоже весят одинаково и меньше настоящих. Как за два взвешивания на весах с
двумя чашами найти 6~фальшивых монет? 

\q2. Дан угол с вершиной $O$, величина которого меньше $180^\circ$, и число $s>0$. На сторонах угла выбирают точки $X$ и $Y$ так, что площадь треугольника $OXY$ равна~$s$. Докажите, что существуют такие точки 
$A$ и $B$ (не~зависящие от выбора точек $X$ и~$Y$), что точки $A$, $B$,
$X$, $Y$ всегда лежат на одной окружности. 

\q3. Трансглактический Рыцарский Турнир проводился по следующим правилам. На шлеме каждого рыцаря светится число~--- количество поединков Турнира, которые этот рыцарь уже выиграл. Поединок разрешен, только если последние цифры на шлемах соперников совпадают. Проигравший в поединке выбывает из Турнира. Историки доподлинно установили, что число участников отличалось от 1\,000\,000 не более чем на 500, а в конце осталось всего 10~рыцарей, между которыми уже было нельзя провести ни одного поединка. Сколько рыцарей участвовало в~Турнире?

\q4. Докажите, что существуют 100 последовательных натуральных чисел, обладающих следующим свойством: каждое из них можно заменить на сумму каких-то двух его различных делителей, больших 1, так, что снова получатся 100 последовательных натуральных чисел (в некотором порядке).

\dotline

\centerline{Олимпиада 2025 года. II тур. 10 класс. Выводная аудитория.}
\smallskip

\q5. Дан остроугольный треугольник $ABC$ с высотой $AD$ и ортоцентром $H$.
Точки $X$ и $Y$ лежат на лучах $AB$ и $AC$ так, что $\angle ABC = \angle
DHX$ и $\angle ACB = \angle DHY$. Докажите, что точка $D$ и~середины
отрезков $BY$, $CX$ и $XY$ лежат на одной окружности. 

\q6. Пространство разбито на единичные кубики, все вершины которых имеют целые координаты. В каждом кубике красным
цветом написано число; все эти числа, кроме конечного количества,~--- нули. 
Для каждой пары соседних кубиков на их общей грани зеленым цветом
написано число, равное модулю разности чисел в этих кубиках. Докажите, что
куб суммы зеленых чисел не меньше чем сумма кубов красных чисел, умноженная на 216. 

\q7. На доске записаны несколько различных слов из букв некоторого алфавита.
Сумма их длин равна $L$. Для каждой упорядоченной пары слов (в том числе,
совпадающих) посчитаем, сколько раз первое слово содержится в качестве
подслова во втором. (Например, слово $aa$ содержится в слове $aaab$ два
раза.) Пусть $K$~--- сумма полученных результатов. Докажите, что
существуют числа $c$ и $\alpha$ (не зависящие ни от размера алфавита, ни
от набора слов) такие, что $K\leq cL^\alpha$, и для всех таких пар $(c,\alpha)$ найдите наименьшее~$\alpha$. 


\clearpage

\centerline{\sc САНКТ-ПЕТЕРБУРГСКАЯ}
\centerline{\sc ОЛИМПИАДА ШКОЛЬНИКОВ 2025 года ПО МАТЕМАТИКЕ}
\centerline{\sc II тур. \ 11 класс.}
\medskip
\hrule
\bigskip

\q1. В пустой Тихий Омут выпустили 100 красных и 99 синих живоглотов.
Когда живоглот хочет есть, он может съесть двух других живоглотов. Если
его жертвы одноцветные, то живоглот меняет свой цвет, а если разноцветные~--- не меняет. В конце остался только один живоглот. Какого он
цвета?

 
\q2. Докажите, что  \
$\displaystyle\frac{1+\tg 1^\circ}{1-\tg 1^\circ}+\frac{1+\tg 2^\circ}{1-\tg 2^\circ}+\ldots+ \frac{1+\tg 15^\circ}{1-\tg 15^\circ}>19$.

\q3. Простое число $p$ обладает свойством:

\smallskip
{\leftskip5mm\rightskip5mm\sl
для всех целых $a$ и $b$ последовательность, заданная условиями 

\centerline{$x_1=a$, $x_2=b$ \ и \ $x_{n+2}=x_{n+1}+x_n$ при $n\geq 1$,}

содержит член, кратный $p$.\par}

\smallskip
Последовательность чисел  Фибоначчи  задаётся соотношениями 

\centerline{$F_1=F_2=1$ и $F_{n+2}=F_{n+1}+F_n$ при $n\geq 1$.} 

Докажите, что $F_{p+1}$ делится на $p$.
 
\q4. На биссектрисе угла с вершиной $A$ зафиксирована точка $B$. На отрезке $AB$ выбирается произвольная точка $O$, для которой окружность с центром $O$ и радиусом $OB$ пересекает стороны угла в четырех
точках. Соединим отрезком две из этих точек, лежащие на разных сторонах
угла и на разных расстояниях от $A$. Докажите, что все построенные таким
образом отрезки касаются одной фиксированной окружности.

\dotline

\centerline{Олимпиада 2025 года. II тур. 11 класс. Выводная аудитория.}
\smallskip

\q5. Дано натуральное число $m$. Докажите, что существует бесконечно
много натуральных $n$, взаимно простых с $m$, для которых
число $(n!)^2+1$~--- составное.
 
\q6. В связном графе больше 1000 вершин и он остаётся связным, даже если из него удалить любые 100 вершин. Докажите, что в~этом графе найдется такой путь $a_1a_2 \ldots a_{102}$, что при удалении этих 102 вершин граф останется связным. (Некоторые пары несоседних вершин пути также могут быть смежны.)
 
\q7. Набор чисел $x_1$, $x_2$,~\dots, $x_{10}$ будем называть \emph{хорошим}, если 

\smallskip
\centerline{$|x_{i+1}-x_1|\geq 3|x_i-x_1|$ \ при $i=1,2,\ldots,9$.} 
\smallskip

Докажите, что из любых 314\,159 вещественных чисел можно выбрать хороший набор.


\end{document}
